% =====================================================================
%  Conferenze, Zagarolo (On The Road), giorno 2 — Talk 5
%  Giochi da tavolo e videogiochi: l'anima digitale di un gioco di
%  carte e l'adattamento di S.T.A.L.K.E.R. in gioco da tavolo
%  Fonte: trascrizione automatica "Track 19". Doppio talk in inglese di
%  due game designer di Awaken Realms, con domande dal pubblico.
%  Etichette con prefisso "c3-5:". Nessun pacchetto aggiuntivo.
% =====================================================================

\chapter[Giochi da tavolo e videogiochi]{Giochi da tavolo e videogiochi: smettere
di fare porting, iniziare ad adattare}
\label{cap:c3-5-giochi-da-tavolo}

L'ultimo intervento della conferenza porta il discorso fuori dallo schermo, sul
\textbf{tavolo}. Due designer di un editore di giochi da tavolo mostrano il
rapporto tra i due media in entrambe le direzioni: prima come un gioco di carte
abbia scoperto la propria ``anima da videogioco'' quando è stato necessario
testarlo con il machine learning; poi come si adatta un videogioco in un gioco da
tavolo senza limitarsi a ``portarlo''. Il capitolo è organizzato così:
\begin{enumerate}
  \item i relatori e il gioco come istinto ancestrale
        (sezione~\ref{sec:c3-5-relatori});
  \item un gioco di carte con un'anima da videogioco: i numeri 7, 5, 4
        (sezione~\ref{sec:c3-5-754});
  \item testare 25 miliardi di combinazioni: l'IA e lo \emph{hate drafting}
        (sezione~\ref{sec:c3-5-test});
  \item porting o adattamento? (sezione~\ref{sec:c3-5-porting});
  \item il caso \emph{S.T.A.L.K.E.R.}: tradurre le sensazioni in meccaniche
        (sezione~\ref{sec:c3-5-stalker});
  \item le lezioni dell'adattamento (sezione~\ref{sec:c3-5-lezioni});
  \item le domande del pubblico (sezione~\ref{sec:c3-5-domande}).
\end{enumerate}

% ---------------------------------------------------------------------
\section{I relatori e il gioco come istinto ancestrale}
\label{sec:c3-5-relatori}

I relatori lavorano entrambi per \textbf{Awaken Realms}\footnote{Trascritto
``related runs'' e ``Away Careers''. Awaken Realms è un editore polacco di giochi
da tavolo noto per le grandi produzioni con miniature.}.
\begin{itemize}
  \item Il primo relatore\footnote{Chiamato ``Gonzalo'' e ``Nuzano'' dal collega;
        il cognome non è pronunciato.} è \textbf{game director} in un'azienda che fa
        giochi da tavolo e videogiochi, con responsabilità soprattutto sui giochi da
        tavolo. In un'industria molto più piccola, il game director fa anche da
        \textbf{production lead}: ogni asset è prodotto da pochissime persone.
  \item Il secondo relatore\footnote{Presentato dal collega come ``Paolo'' (e poi
        ``Paola'', ``Pablo''), si presenta da sé come ``Otto Zamborski'': entrambe
        le forme sono incerte.} è polacco e lavora con Awaken Realms da circa
        \textbf{dieci anni}; ha curato l'adattamento da tavolo di
        \textbf{\emph{S.T.A.L.K.E.R.}}.
\end{itemize}
Il primo relatore, per sua stessa ammissione, parla ``venti minuti per presentare
un altro''.

\subsection{Il gioco è più antico della cultura}

Il primo relatore apre con \textbf{Johan Huizinga}, lo storico olandese autore di
\emph{Homo Ludens}: \textbf{il gioco è più antico della cultura}. Racconta dei suoi
due gattini: ancora non sanno bere bene, ma a un certo punto si guardano, saltano e
cominciano a giocare, mordendosi ma non troppo, perché l'altro segnala ``mi fai
male''. Anche gli animali giocano: il gioco è qualcosa di \textbf{ancestrale} che
portiamo dentro dalla nascita. E il game design analogico lavora proprio sulle
\textbf{basi}, sui piccoli frammenti di questa esperienza.

\subsection{Un'industria più piccola, ma che cresce}

Pochi in sala hanno giocato a giochi da tavolo moderni (``\emph{Monopoly} non
conta'', da \emph{Catan} in poi). L'industria dei giochi da tavolo è
\textbf{un decimo o meno} di quella dei videogiochi, ma \textbf{crescere} al suo
interno è molto più rapido e facile. Chi passa al design analogico mantiene la
stessa esperienza da game designer, ma trova un ``sandbox'' in cui crescere; un
argomento non da poco, per una platea preoccupata per il lavoro.

% ---------------------------------------------------------------------
\section{Un gioco di carte con un'anima da videogioco}
\label{sec:c3-5-754}

\subsection{I vincoli: 7, 5, 4}

Nei giochi da tavolo non ci sono effetti sonori, e poca narrativa a meno di
leggerla. Si lavora quindi sui \textbf{confini}: comprimere l'esperienza perché il
cervello la capisca facilmente e vi si appassioni, senza una ``supermacchina'' che
calcoli tutto dietro le quinte. Per il loro gioco i vincoli sono tre numeri, ispirati
dalla passione per \textbf{\emph{Magic: The Gathering}}\footnote{Trascritto
``Mother of the Magic the Gallery''.}:
\begin{description}
  \item[7] carte in draft per giocatore;
  \item[5] dadi sul tavolo, che sono il \textbf{mana}, cioè la valuta;
  \item[4] azioni possibili: a ogni turno si sceglie una sola azione tra quattro.
\end{description}
Il gioco che ne è nato\footnote{Il titolo è trascritto ``Either Stone''; non è
ricostruibile con sicurezza.} è un fantascientifico ``soft'' in cui si è
un'entità che combatte il \textbf{vuoto}, metafora della vita che smette di
esistere: si lotta contro qualcosa di inevitabile, ma si lotta comunque.

\subsection{Lo spazio sulla carta e il regolamento}

Nei giochi da tavolo la \textbf{direzione artistica} conta moltissimo: di solito
l'illustrazione occupa \textbf{due terzi} della carta, e nello spazio restante si
può comunicare pochissimo. La disposizione:
\begin{itemize}
  \item il \textbf{costo} in alto a sinistra;
  \item i \textbf{punti vittoria} sul lato opposto;
  \item sotto, l'\textbf{attivazione} e le \textbf{abilità}.
\end{itemize}
Un regolamento dovrebbe essere \textbf{minuscolo}: leggere un regolamento non è
ancora giocare, e non c'è una demo che insegni. Per questo hanno scelto un
\textbf{regolamento essenziale} e messo tutte le altre regole \textbf{sulle carte}.

\subsection{Le azioni e il motore di combo}

\begin{itemize}
  \item \textbf{Prendere un dado}: è un \emph{open draft}. Si tirano i cinque dadi,
        se ne prende uno e si ottengono due mana di quel colore.
  \item \textbf{Giocare una carta}: le carte riportano il numero o il colore del
        dado appena preso; a quel punto il dado ``salta'' da una carta all'altra
        innescando una \textbf{reazione a catena} di tutto ciò che si è pianificato.
        Ma con sole sette carte il motore non è un mazzo da 40--60 come in
        \emph{Magic}: se ne giochi una, ne restano sei.
  \item \textbf{Attaccare} e \textbf{riposare} (rimettere in piedi le truppe).
\end{itemize}
L'open draft è la parte più interessante, quella che evita al gioco di sembrare una
copia dei suoi modelli.

\subsection{La curva nascosta e l'economia}

Testando, hanno costruito una \textbf{curva nascosta}: pagando 2 mana si ottiene 1
punto vittoria, pagandone 8 se ne ottengono 4. Facile. La parte difficile è
l'\textbf{incastro delle abilità}: volevano che una singola azione potesse
diventarne sei o sette nello stesso turno. Siccome il mana non bastava, hanno
introdotto la \textbf{vita come seconda valuta}. Tutto questo, nota il relatore,
somiglia molto allo sviluppo dell'ecosistema di un videogioco: \textbf{valute,
abilità, trigger}.

\subsection{Tre condizioni di fine partita}

La partita finisce in tre modi:
\begin{enumerate}
  \item un giocatore ha giocato tutte e sette le carte;
  \item tutte le minacce sono state sconfitte;
  \item i \textbf{punti vittoria finiscono}: in una scatola non si possono avere
        punti infiniti, ogni cosa deve avere un limite.
\end{enumerate}

% ---------------------------------------------------------------------
\section{Testare 25 miliardi di combinazioni}
\label{sec:c3-5-test}

\subsection{Il problema dei numeri}

Combinando le sette carte che si possono ottenere si arriva a quasi \textbf{due
miliardi} di combinazioni. Aggiungendo i \textbf{leader}, ciascuno con abilità e
trigger propri, e moltiplicando per le \textbf{75 carte} del gioco base e i
\textbf{13 leader}, si arriva a \textbf{25 miliardi}. Nessun essere umano può
verificare che nulla sia rotto. Testando al tavolo con tutto il team, dopo 40 minuti
avevano completato \textbf{un solo test}.

\subsection{L'IA che non gioca come un umano}

Qui il gioco da tavolo ``abbraccia il videogioco'': hanno usato il
\textbf{machine learning} per testare tutte le combinazioni in pochi secondi. Ma è
emerso un problema. Nelle partite umane la fine per ``tutte le carte giocate''
avveniva in circa \textbf{un terzo} dei casi; con l'IA nel \textbf{doppio} dei casi.
L'IA giocava carte velocissimamente, mentre lo scopo del gioco è usare le abilità
per fare punti, non giocare carte a ogni costo: \textbf{non capiva come si comporta
un umano}.

Hanno aggiunto rami di logica sul draft, ma l'IA non vedeva una potenzialità della
mente umana: lo \textbf{hate drafting}, cioè prendere un dado non perché serve a sé
ma per togliere punti vittoria a un avversario. È un comportamento così vario che
spiegarlo all'IA era impossibile: l'IA doveva \textbf{imparare a giocare come un
umano}.

\subsection{La versione digitale}

Così hanno fatto il \textbf{videogioco} (``scusate i colleghi in fondo alla sala: non
faccio videogiochi per mestiere, vi sembrerà un prototipo''). Hanno scoperto che il
gioco da tavolo si poteva giocare anche come videogioco: la sua \textbf{anima da
videogioco}, il legame tra i due temi del talk. Invitando persone a giocare, negli
stessi 40 minuti sono passati da uno a \textbf{80 test}.

L'espansione va consegnata a febbraio dell'anno successivo e il tempo stringe, ma
lo strumento ha fatto crescere i contenuti oltre le attese: volevano testare 70
carte, ora sono circa \textbf{200}. Aiuta a testare e ad anticipare che cosa faranno
i giocatori. Poi ci saranno comunque i test umani, perché bisogna trovare la
\textbf{scintilla}, ``come diceva Jonas ieri''\footnote{Si veda il
capitolo~\ref{cap:c2-5-raw-fury}.}: deve essere divertente.

La versione digitale è giocabile online, nel browser (``tutto JavaScript''), con una
demo che sostituisce il regolamento. Per evitare che il test fosse noioso
(``abbiamo giocato, e adesso? Dammi qualcosa''), l'hanno \textbf{gamificata}:
obiettivi giornalieri, taglie, esperienza per il personaggio. Per il relatore questo
apre un mondo di possibilità nella creazione di giochi da tavolo.

% ---------------------------------------------------------------------
\section{Porting o adattamento?}
\label{sec:c3-5-porting}

Il secondo talk si intitola, più o meno, ``\textbf{smettete di fare porting dei
videogiochi e iniziate a reinventarli}''. Oggi è comune prendere una IP di
videogiochi e farne un gioco da tavolo, e spesso, ``mi dispiace dirlo'', è
un'operazione puramente commerciale. Il primo relatore lo introduce così: il gioco
da tavolo ha un'anima da videogioco e il porting sembra naturale, ma non lo è; un
port ``uno a uno'' riesce male e non diverte quanto ci si aspetta. Bisogna
\textbf{sottrarre la scintilla}, le emozioni del videogioco, e portarle in
un'esperienza fisica.

\subsection{Due media diversi}

Videogiochi e giochi da tavolo sono \textbf{media diversi}, come un romanzo e un
film. Condividono la stessa anima e possono suscitare emozioni simili, ma con
\textbf{strumenti diversi}.
\begin{description}
  \item[Porting] è la stessa cosa su un hardware diverso: come tradurre un libro in
        un'altra lingua. Alla fine si ha lo stesso libro.
  \item[Adattamento] è riconoscere che il medium è diverso: trovare ciò che rende
        quel gioco \emph{quel} gioco, e prendere decisioni che gli autori originali
        non hanno dovuto prendere. Molte cose si traducono uno a uno, ma
        \textbf{non sempre dovrebbero}.
\end{description}

\subsection{I giochi da tavolo girano sugli umani}

La differenza fondamentale non è schermo contro cartone (a volte i giochi da tavolo
finiscono persino dentro i videogiochi, come in uno screenshot di
\emph{StarCraft II}). È che i videogiochi offrono \textbf{scelte} e mettono alla
prova l'\textbf{abilità manuale}: in \emph{S.T.A.L.K.E.R.}, FPS post-apocalittico
ambientato nell'Ucraina contemporanea intorno a Chernobyl, non conta solo
\emph{se} sparo ma \emph{come} sparo. Nei giochi da tavolo di solito c'è solo la
\textbf{scelta}: mira e riflessi, che nei videogiochi si danno per scontati, non
esistono o si traducono di rado.

Inoltre il gioco gira su un \textbf{hardware non standardizzato}: le persone. Si
può assumere che abbiano due mani, ma per il resto hanno background diversi e
capiscono le regole in modo diverso.

\subsection{Limiti e punti di forza del medium}

Limiti:
\begin{itemize}
  \item \textbf{abilità manuale}, come detto;
  \item \textbf{effetti passivi}: un computer li calcola gratis in background; al
        tavolo le persone \textbf{se li dimenticano}, a meno che non siano
        \textbf{positivi}, che restano più impressi;
  \item \textbf{calcoli pesanti}: chi gioca da tavolo vuole divertirsi, non fare
        matematica.
\end{itemize}
Tutto si può aggirare, ma con risultati variabili.

Punti di forza:
\begin{itemize}
  \item la ``\textbf{risoluzione}'' dell'interfaccia: un gioco da tavolo mostra
        moltissime informazioni insieme, con tutte le carte disposte sul tavolo; il
        problema è solo se il giocatore saprà decidere con tutto ciò che ha davanti.
        Nel porting inverso, che il relatore ha fatto una volta, una mano di carte
        che al tavolo era sempre visibile era nascosta nel formato digitale, e ci è
        voluto molto per aggirare il problema;
  \item l'\textbf{interazione tra giocatori}: online è molto ridotta; al tavolo ci
        si vede e si interagisce anche in modo non verbale. Lo hate drafting fatto
        alla persona seduta accanto è tutt'altra cosa: le battute, il divertimento
        sociale. Le regole sono spesso solo un pretesto per divertirsi insieme.
\end{itemize}

% ---------------------------------------------------------------------
\section{Il caso \emph{S.T.A.L.K.E.R.}}
\label{sec:c3-5-stalker}

\subsection{Trovare ciò che fa brillare il gioco}

Il primo passo, ``come diceva Jonas'', è trovare \textbf{ciò che fa brillare} il
gioco. A volte è solo un FPS qualunque degli anni 2010; nel caso di
\emph{S.T.A.L.K.E.R.} bisognava diagnosticare che cosa lo rende sé stesso.
\begin{itemize}
  \item L'\textbf{atmosfera} era facile: si traduce quasi senza perdite, grazie
        alla grafica e al tono visivo (anche se non ci sarà una colonna sonora). E
        tutto il team era fan della serie: non c'era bisogno di spiegare che cosa
        fare.
  \item Difficile era il \textbf{mondo aperto}: far stare un gioco da 50 ore in
        sessioni da circa \textbf{quattro ore}, e rendere ogni sessione abbastanza
        interessante da far tornare le persone per la successiva.
\end{itemize}

\subsection{Il ritorno del combattimento}

All'inizio il relatore non voleva fare l'ennesimo dungeon crawler ed
\textbf{evitava il combattimento}, perché è pesante di meccaniche, quindi meno
accessibile. Poi arrivò il ``bagno di realtà'' dai giocatori: ``state adattando un
FPS: dov'è la sparatoria?''. Si tornò al tavolo da disegno e il combattimento rientrò,
con tutte le miniature. Ma questo aggiunge \textbf{regole}, e gli stakeholder devono
accettare che non sarà un gioco facile da imparare.

\subsection{Dalle sensazioni alle meccaniche}

Trovata la scintilla, si individuano le \textbf{sensazioni} che il gioco suscita e
si cerca come tradurle con le meccaniche che funzionano meglio nel nuovo medium.

\paragraph{Le anomalie.} Nella Zona ci sono luoghi in cui la fisica non funziona:
nel videogioco sono trappole più o meno visibili, basate sulla destrezza; se il
giocatore non le vede, muore. In un gioco da tavolo morire significa
\textbf{ricominciare}, con enorme perdita di tempo. Al tavolo le anomalie sono
diventate un sistema di \textbf{incertezza}: gli spazi hanno valori diversi di
probabilità di innesco; si può entrare alla cieca sperando nella fortuna, oppure
lanciare \textbf{bulloni e dadi} (come nel videogioco) per annullare l'effetto. La
sensazione è la stessa: rischiare, oppure spendere risorse per andare sul sicuro.

\paragraph{Gli artefatti.} Sono i mattoni del personaggio, con abilità diverse. Nel
videogioco un artefatto dà, per esempio, ``+5\% di resistenza ai proiettili'',
``che nessuno sa cosa significhi''. Al tavolo non si vuole calcolare il 5\%, e
\textbf{niente passivi}: ogni effetto dev'essere un'\textbf{azione}. L'artefatto
diventa: prima del tiro di difesa, puoi tirare un dado e ridurre il danno. Una
scelta del giocatore invece di un effetto passivo.

\paragraph{Lo sparo.} Nei giochi dei propri ``anni formativi'' si prende in
prestito ciò che piaceva da giocatore. In un FPS si sceglie dove mirare; al tavolo
si sceglie tra \textbf{colpo alla testa} e \textbf{colpo al corpo}, una meccanica
presa da un gioco che il relatore considera ``uno dei migliori mai fatti''\footnote{Il
titolo è irriconoscibile nella trascrizione (``the first product'', ``Pauly'').}.

\paragraph{I componenti.} I giochi da tavolo sono \textbf{oggetti fisici}: quando si
può, conviene usare componenti che facciano sentire il mondo. I segnalini dei
bulloni hanno la forma di bulloni; i segnalini munizioni, usati per ricaricare, hanno
la forma di \textbf{caricatori}. È una sensazione tattile di ciò che si sta facendo.

\paragraph{L'IA dei nemici.} È stato lo ``stress test'' di quanto possa reggere
l'hardware umano. \emph{S.T.A.L.K.E.R.} si basa su fazioni in guerra tra loro, e i
giocatori se lo aspettavano anche al tavolo; il primo capitolo era famoso per il
suo sistema di IA che simulava la vita della Zona. Al tavolo l'IA doveva essere
piuttosto pesante: inizialmente somigliava ai blueprint di Unreal Engine, poi è
stata semplificata in \textbf{carte con istruzioni}. Ma un sistema complesso
significa molta \textbf{gestione al tavolo}: ogni istruzione va eseguita da un
giocatore, e quando un giocatore non fa scelte \textbf{non sta giocando}, ``fa il
lavoro del computer''. Bisogna tagliare il più possibile, mantenendo un'IA
\textbf{intelligente e sorprendente}.

\paragraph{Il rumore.} \emph{S.T.A.L.K.E.R.} ha qualcosa del genere \emph{immersive
sim}: i problemi si possono affrontare in modi diversi. Il relatore ha creato una
meccanica di \textbf{rumore}, un sistema di \emph{aggro}: i nemici seguono il
giocatore quando lo sentono o lo vedono. L'ispirazione viene da un gioco
storico\footnote{Non nominato in modo riconoscibile; la descrizione (spazi che
generano rumore, guardie che lo seguono) fa pensare a \emph{Thief: The Dark
Project}, ma è un'ipotesi.} in cui molte superfici generavano rumore e le guardie lo
seguivano. Il relatore si ispira spesso ai giochi degli \textbf{anni Ottanta e
Novanta}, che per i limiti di memoria avevano vincoli simili a quelli dei giochi da
tavolo. Il risultato: i giocatori hanno imparato che il \textbf{rumore è una
risorsa} (fai rumore, attiri i nemici, li porti nelle trappole), ed è nato un
\textbf{gioco stealth} da tavolo, genere raro, che ha reso l'adattamento un gioco
unico.

% ---------------------------------------------------------------------
\section{Le lezioni dell'adattamento}
\label{sec:c3-5-lezioni}

\begin{enumerate}
  \item \textbf{Una sensazione non è un insieme di meccaniche.} Trovare che cosa fa
        brillare il gioco di partenza, poi capire come dovrebbe funzionare nel medium
        di arrivo.
  \item \textbf{Decidere se si adatta la storia o l'esperienza}: si otterranno giochi
        diversi.
  \item \textbf{Conoscere la fonte, ma non esserne prigionieri.} Il relatore ama la
        serie con tutto il cuore, ma questo lo ha reso ``sordo'' al feedback di chi
        non la conosceva, e il gioco è diventato più complicato del necessario.
  \item \textbf{Giocare il gioco due volte}: prima da giocatore, per capire che cosa
        sentono le persone; poi da designer, per capire che cosa pensavano gli
        autori quando hanno preso le loro decisioni.
\end{enumerate}
``A volte \textbf{gli adattamenti peggiori sono i più fedeli}.'' Ha qualche esempio,
ma preferisce non farlo in pubblico.

% ---------------------------------------------------------------------
\section{Le domande del pubblico}
\label{sec:c3-5-domande}

\subsection{Quanto costa prototipare un gioco da tavolo?}

Per chi fa videogiochi prototipare ``non costa nulla''; e per un gioco da tavolo?
\begin{itemize}
  \item Il prototipo di base costa poco: carta, dadi e carte di recupero, una
        stampante.
  \item Awaken Realms testa i giochi \textbf{in digitale} fino al momento in cui
        bisogna verificare UI e UX fisiche: allora il gioco va sul tavolo, per
        capire se i movimenti delle mani sono scomodi o troppo macchinosi.
  \item Per il primo gioco in assoluto, meglio \textbf{evitare le miniature}: sono
        il costo maggiore. Uno \textbf{stampo} dura moltissimo, ma costa circa
        \textbf{5000--6000 euro}; Awaken Realms ha a volte centinaia di miniature per
        gioco, e quelle grandi non entrano in un solo stampo: vanno divise in pezzi
        e preassemblate.
  \item Esistono stampanti 3D per avere il pezzo sul tavolo, ma un file STL può avere
        punti fragili che si rompono durante il trasporto oltreoceano.
  \item Ci sono servizi di stampa su richiesta\footnote{Il servizio è trascritto
        ``Lounge Tabletop'' e non è ricostruibile con sicurezza.}: si può ordinare
        anche una sola copia, a un costo unitario molto più alto.
\end{itemize}

\subsection{Giochi da tavolo completamente digitalizzati?}

Esistono siti che digitalizzano interamente i giochi da tavolo per giocare in
multiplayer. Funziona (mescolare un mazzo richiede un clic), ma non si perde la
sensazione del tavolo? Per il primo relatore dipende dall'esperienza:
\begin{itemize}
  \item nei \textbf{giochi sociali} e di deduzione, dove i giocatori stessi sono i
        ``componenti'' che rendono il gioco divertente, il fattore umano è il cuore
        e la digitalizzazione perde molto;
  \item molti \textbf{german} (\emph{eurogame}) sono invece \textbf{migliori in
        digitale}, perché tutta la gestione è fatta dallo script; nei
        \textbf{tornei} il digitale elimina anche il fattore umano e i problemi come
        le carte segnate;
  \item \textbf{Board Game Arena}\footnote{Trascritto ``VGA'' e ``BGA'',
        ``boarding''.} è un ottimo canale pubblicitario: molti provano il gioco lì e
        poi comprano la copia fisica; ha inoltre un modello ``alla Spotify'' che
        remunera in base alle partite.
\end{itemize}

\subsection{I giochi ibridi}

Un partecipante cita \emph{Mansions of Madness}\footnote{Trascritto ``Machine of
Themis'' e ``mention of madness''.}, gioco da tavolo con un'app. Per i relatori gli
ibridi vanno bene \textbf{per chi gioca ai videogiochi, non per chi gioca da
tavolo}. Il pubblico è un \textbf{imbuto}: si comincia con \emph{Monopoly}, poi
\emph{Exploding Kittens}, e magari ci si ferma lì per sempre; chi vuole più
stimoli passa a giochi più complessi e diventa esperto. Un esperto che gioca anche
ai videogiochi capisce l'ibrido; chi viene solo dal tavolo difficilmente lo
accetta. Awaken Realms ha chiesto più volte alla propria community se volesse un
gioco ibrido, trovando forte avversione. Resta però una sfida interessante per un
video game designer: per esempio un gioco sociale in cui delle voci in cuffia
sussurrano a ciascuno che cosa fare, con prove fisiche da muovere sul tavolo e il
digitale che gestisce i compiti ripetitivi.

\subsection{Il caso \emph{XCOM}}

Un partecipante cita il gioco da tavolo di \emph{XCOM}, che quasi rinuncia alla parte
tattica. Il secondo relatore spiega: ha tagliato la mappa tattica ed è diventato un
gestionale \textbf{in tempo reale} guidato da un'app, cioè fa una cosa che i
giocatori da tavolo di solito odiano: vogliono \textbf{tempo per pensare}. Il tempo
reale va bene nei party game. Lui avrebbe affrontato la cosa diversamente, anche se
il gioco ``funzionava abbastanza bene''. In generale bisogna \textbf{tagliare
qualcosa} dall'originale: in \emph{S.T.A.L.K.E.R.} ha messo quasi tutto, ma ha
tagliato gran parte dell'esplorazione della Zona, sostituendola con una mappa e
degli adesivi per le basi.

% ---------------------------------------------------------------------
\section*{Punti chiave del capitolo}
\addcontentsline{toc}{section}{Punti chiave del capitolo}

\begin{itemize}
  \item Il gioco è più antico della cultura (Huizinga); il design analogico lavora
        sulle basi dell'esperienza, con vincoli forti (7 carte, 5 dadi, 4 azioni).
  \item Regolamento minimo, regole sulle carte; valute, abilità e trigger come in un
        ecosistema da videogioco.
  \item Con miliardi di combinazioni serve il test automatico, ma l'IA non gioca come
        un umano (\emph{hate drafting}): una versione digitale gamificata ha
        moltiplicato i test.
  \item Porting e adattamento sono cose diverse: i giochi da tavolo girano sugli
        umani, non gestiscono bene abilità manuale, passivi e calcoli, ma mostrano
        più informazioni e hanno più interazione sociale.
  \item Adattare \emph{S.T.A.L.K.E.R.}: anomalie come scelta di rischio, artefatti
        come azioni, colpo alla testa o al corpo, componenti evocativi, IA a carte,
        rumore come risorsa.
  \item Una sensazione non è un insieme di meccaniche; conoscere la fonte senza
        esserne prigionieri; gli adattamenti peggiori sono spesso i più fedeli.
  \item Prototipi economici, miniature costose; digitale utile per eurogame e
        marketing, meno per i giochi sociali; ibridi graditi ai videogiocatori più
        che ai giocatori da tavolo.
\end{itemize}
